% Options for packages loaded elsewhere
\PassOptionsToPackage{unicode}{hyperref}
\PassOptionsToPackage{hyphens}{url}
%
\documentclass[
]{article}
\usepackage{amsmath,amssymb}
\usepackage{iftex}
\ifPDFTeX
  \usepackage[T1]{fontenc}
  \usepackage[utf8]{inputenc}
  \usepackage{textcomp} % provide euro and other symbols
\else % if luatex or xetex
  \usepackage{unicode-math} % this also loads fontspec
  \defaultfontfeatures{Scale=MatchLowercase}
  \defaultfontfeatures[\rmfamily]{Ligatures=TeX,Scale=1}
\fi
\usepackage{lmodern}
\ifPDFTeX\else
  % xetex/luatex font selection
\fi
% Use upquote if available, for straight quotes in verbatim environments
\IfFileExists{upquote.sty}{\usepackage{upquote}}{}
\IfFileExists{microtype.sty}{% use microtype if available
  \usepackage[]{microtype}
  \UseMicrotypeSet[protrusion]{basicmath} % disable protrusion for tt fonts
}{}
\makeatletter
\@ifundefined{KOMAClassName}{% if non-KOMA class
  \IfFileExists{parskip.sty}{%
    \usepackage{parskip}
  }{% else
    \setlength{\parindent}{0pt}
    \setlength{\parskip}{6pt plus 2pt minus 1pt}}
}{% if KOMA class
  \KOMAoptions{parskip=half}}
\makeatother
\usepackage{xcolor}
\usepackage[margin=1in]{geometry}
\usepackage{color}
\usepackage{fancyvrb}
\newcommand{\VerbBar}{|}
\newcommand{\VERB}{\Verb[commandchars=\\\{\}]}
\DefineVerbatimEnvironment{Highlighting}{Verbatim}{commandchars=\\\{\}}
% Add ',fontsize=\small' for more characters per line
\usepackage{framed}
\definecolor{shadecolor}{RGB}{248,248,248}
\newenvironment{Shaded}{\begin{snugshade}}{\end{snugshade}}
\newcommand{\AlertTok}[1]{\textcolor[rgb]{0.94,0.16,0.16}{#1}}
\newcommand{\AnnotationTok}[1]{\textcolor[rgb]{0.56,0.35,0.01}{\textbf{\textit{#1}}}}
\newcommand{\AttributeTok}[1]{\textcolor[rgb]{0.13,0.29,0.53}{#1}}
\newcommand{\BaseNTok}[1]{\textcolor[rgb]{0.00,0.00,0.81}{#1}}
\newcommand{\BuiltInTok}[1]{#1}
\newcommand{\CharTok}[1]{\textcolor[rgb]{0.31,0.60,0.02}{#1}}
\newcommand{\CommentTok}[1]{\textcolor[rgb]{0.56,0.35,0.01}{\textit{#1}}}
\newcommand{\CommentVarTok}[1]{\textcolor[rgb]{0.56,0.35,0.01}{\textbf{\textit{#1}}}}
\newcommand{\ConstantTok}[1]{\textcolor[rgb]{0.56,0.35,0.01}{#1}}
\newcommand{\ControlFlowTok}[1]{\textcolor[rgb]{0.13,0.29,0.53}{\textbf{#1}}}
\newcommand{\DataTypeTok}[1]{\textcolor[rgb]{0.13,0.29,0.53}{#1}}
\newcommand{\DecValTok}[1]{\textcolor[rgb]{0.00,0.00,0.81}{#1}}
\newcommand{\DocumentationTok}[1]{\textcolor[rgb]{0.56,0.35,0.01}{\textbf{\textit{#1}}}}
\newcommand{\ErrorTok}[1]{\textcolor[rgb]{0.64,0.00,0.00}{\textbf{#1}}}
\newcommand{\ExtensionTok}[1]{#1}
\newcommand{\FloatTok}[1]{\textcolor[rgb]{0.00,0.00,0.81}{#1}}
\newcommand{\FunctionTok}[1]{\textcolor[rgb]{0.13,0.29,0.53}{\textbf{#1}}}
\newcommand{\ImportTok}[1]{#1}
\newcommand{\InformationTok}[1]{\textcolor[rgb]{0.56,0.35,0.01}{\textbf{\textit{#1}}}}
\newcommand{\KeywordTok}[1]{\textcolor[rgb]{0.13,0.29,0.53}{\textbf{#1}}}
\newcommand{\NormalTok}[1]{#1}
\newcommand{\OperatorTok}[1]{\textcolor[rgb]{0.81,0.36,0.00}{\textbf{#1}}}
\newcommand{\OtherTok}[1]{\textcolor[rgb]{0.56,0.35,0.01}{#1}}
\newcommand{\PreprocessorTok}[1]{\textcolor[rgb]{0.56,0.35,0.01}{\textit{#1}}}
\newcommand{\RegionMarkerTok}[1]{#1}
\newcommand{\SpecialCharTok}[1]{\textcolor[rgb]{0.81,0.36,0.00}{\textbf{#1}}}
\newcommand{\SpecialStringTok}[1]{\textcolor[rgb]{0.31,0.60,0.02}{#1}}
\newcommand{\StringTok}[1]{\textcolor[rgb]{0.31,0.60,0.02}{#1}}
\newcommand{\VariableTok}[1]{\textcolor[rgb]{0.00,0.00,0.00}{#1}}
\newcommand{\VerbatimStringTok}[1]{\textcolor[rgb]{0.31,0.60,0.02}{#1}}
\newcommand{\WarningTok}[1]{\textcolor[rgb]{0.56,0.35,0.01}{\textbf{\textit{#1}}}}
\usepackage{graphicx}
\makeatletter
\newsavebox\pandoc@box
\newcommand*\pandocbounded[1]{% scales image to fit in text height/width
  \sbox\pandoc@box{#1}%
  \Gscale@div\@tempa{\textheight}{\dimexpr\ht\pandoc@box+\dp\pandoc@box\relax}%
  \Gscale@div\@tempb{\linewidth}{\wd\pandoc@box}%
  \ifdim\@tempb\p@<\@tempa\p@\let\@tempa\@tempb\fi% select the smaller of both
  \ifdim\@tempa\p@<\p@\scalebox{\@tempa}{\usebox\pandoc@box}%
  \else\usebox{\pandoc@box}%
  \fi%
}
% Set default figure placement to htbp
\def\fps@figure{htbp}
\makeatother
\setlength{\emergencystretch}{3em} % prevent overfull lines
\providecommand{\tightlist}{%
  \setlength{\itemsep}{0pt}\setlength{\parskip}{0pt}}
\setcounter{secnumdepth}{-\maxdimen} % remove section numbering
\usepackage{graphicx}
\usepackage{bookmark}
\IfFileExists{xurl.sty}{\usepackage{xurl}}{} % add URL line breaks if available
\urlstyle{same}
\hypersetup{
  hidelinks,
  pdfcreator={LaTeX via pandoc}}

\author{}
\date{\vspace{-2.5em}}

\begin{document}

\begin{Shaded}
\begin{Highlighting}[]
\FunctionTok{options}\NormalTok{(}\AttributeTok{OutDec =} \StringTok{"."}\NormalTok{)  }\CommentTok{\# Asegurar punto decimal en este chunk}

\CommentTok{\# Establecer semilla aleatoria para reproducibilidad}
\FunctionTok{set.seed}\NormalTok{(}\FunctionTok{sample}\NormalTok{(}\DecValTok{1}\SpecialCharTok{:}\DecValTok{10000}\NormalTok{, }\DecValTok{1}\NormalTok{))}

\CommentTok{\# Aleatorizar contexto empresarial}
\NormalTok{tipos\_empresa }\OtherTok{\textless{}{-}} \FunctionTok{c}\NormalTok{(}\StringTok{"empresa"}\NormalTok{, }\StringTok{"compañía"}\NormalTok{, }\StringTok{"organización"}\NormalTok{, }\StringTok{"corporación"}\NormalTok{, }\StringTok{"entidad"}\NormalTok{, }\StringTok{"firma"}\NormalTok{, }\StringTok{"institución"}\NormalTok{, }\StringTok{"negocio"}\NormalTok{)}
\NormalTok{tipo\_empresa }\OtherTok{\textless{}{-}} \FunctionTok{sample}\NormalTok{(tipos\_empresa, }\DecValTok{1}\NormalTok{)}

\NormalTok{sectores\_empresa }\OtherTok{\textless{}{-}} \FunctionTok{c}\NormalTok{(}\StringTok{"tecnológica"}\NormalTok{, }\StringTok{"financiera"}\NormalTok{, }\StringTok{"industrial"}\NormalTok{, }\StringTok{"comercial"}\NormalTok{, }\StringTok{"de servicios"}\NormalTok{, }\StringTok{"consultora"}\NormalTok{, }\StringTok{"energética"}\NormalTok{,}
                     \StringTok{"automotriz"}\NormalTok{, }\StringTok{"farmacéutica"}\NormalTok{, }\StringTok{"alimentaria"}\NormalTok{, }\StringTok{"textil"}\NormalTok{, }\StringTok{"educativa"}\NormalTok{)}
\NormalTok{sector\_empresa }\OtherTok{\textless{}{-}} \FunctionTok{sample}\NormalTok{(sectores\_empresa, }\DecValTok{1}\NormalTok{)}

\NormalTok{terminos\_encuesta }\OtherTok{\textless{}{-}} \FunctionTok{c}\NormalTok{(}\StringTok{"encuesta"}\NormalTok{, }\StringTok{"sondeo"}\NormalTok{, }\StringTok{"estudio"}\NormalTok{, }\StringTok{"investigación"}\NormalTok{, }\StringTok{"censo"}\NormalTok{, }\StringTok{"análisis"}\NormalTok{, }\StringTok{"relevamiento"}\NormalTok{, }\StringTok{"consulta"}\NormalTok{)}
\NormalTok{termino\_encuesta }\OtherTok{\textless{}{-}} \FunctionTok{sample}\NormalTok{(terminos\_encuesta, }\DecValTok{1}\NormalTok{)}

\CommentTok{\# Aleatorizar bienes}
\NormalTok{bienes\_vivienda }\OtherTok{\textless{}{-}} \FunctionTok{c}\NormalTok{(}
  \StringTok{"casa"}\NormalTok{, }\StringTok{"apartamento"}\NormalTok{, }\StringTok{"vivienda"}\NormalTok{, }\StringTok{"piso"}\NormalTok{, }\StringTok{"chalet"}\NormalTok{, }\StringTok{"dúplex"}\NormalTok{, }\StringTok{"residencia"}\NormalTok{, }\StringTok{"inmueble"}
\NormalTok{)}
\NormalTok{bien\_vivienda1 }\OtherTok{\textless{}{-}} \FunctionTok{sample}\NormalTok{(bienes\_vivienda, }\DecValTok{1}\NormalTok{)}
\NormalTok{bienes\_vivienda\_restantes }\OtherTok{\textless{}{-}}\NormalTok{ bienes\_vivienda[bienes\_vivienda }\SpecialCharTok{!=}\NormalTok{ bien\_vivienda1]}
\NormalTok{bien\_vivienda2 }\OtherTok{\textless{}{-}} \FunctionTok{sample}\NormalTok{(bienes\_vivienda\_restantes, }\DecValTok{1}\NormalTok{)}

\NormalTok{bienes\_transporte }\OtherTok{\textless{}{-}} \FunctionTok{c}\NormalTok{(}
  \StringTok{"carro"}\NormalTok{, }\StringTok{"automóvil"}\NormalTok{, }\StringTok{"vehículo"}\NormalTok{, }\StringTok{"coche"}
\NormalTok{)}
\NormalTok{bien\_transporte }\OtherTok{\textless{}{-}} \FunctionTok{sample}\NormalTok{(bienes\_transporte, }\DecValTok{1}\NormalTok{)}

\CommentTok{\# Aleatorizar los porcentajes del gráfico circular manteniendo coherencia}
\CommentTok{\# Valores originales: 24\% (carro y apartamento), 20\% (solo apartamento), 31\% (solo casa),}
\CommentTok{\#                     6\% (solo carro), 19\% (carro y casa)}

\CommentTok{\# Generar variaciones de los porcentajes manteniendo suma = 100\% y relaciones coherentes}
\NormalTok{total\_puntos\_porcentaje }\OtherTok{\textless{}{-}} \DecValTok{100}

\CommentTok{\# Aleatorizar el porcentaje para "bien\_transporte y bien\_vivienda1" (equivalente a "carro y apartamento")}
\NormalTok{p\_bien\_transporte\_y\_vivienda1 }\OtherTok{\textless{}{-}} \FunctionTok{sample}\NormalTok{(}\DecValTok{18}\SpecialCharTok{:}\DecValTok{28}\NormalTok{, }\DecValTok{1}\NormalTok{)}

\CommentTok{\# Aleatorizar el porcentaje para "solo bien\_vivienda2" (equivalente a "solo casa")}
\NormalTok{p\_solo\_vivienda2 }\OtherTok{\textless{}{-}} \FunctionTok{sample}\NormalTok{(}\DecValTok{25}\SpecialCharTok{:}\DecValTok{35}\NormalTok{, }\DecValTok{1}\NormalTok{)}

\CommentTok{\# Aleatorizar el porcentaje para "bien\_transporte y bien\_vivienda2" (equivalente a "carro y casa")}
\NormalTok{p\_bien\_transporte\_y\_vivienda2 }\OtherTok{\textless{}{-}} \FunctionTok{sample}\NormalTok{(}\DecValTok{15}\SpecialCharTok{:}\DecValTok{25}\NormalTok{, }\DecValTok{1}\NormalTok{)}

\CommentTok{\# Aleatorizar el porcentaje para "solo bien\_transporte" (equivalente a "solo carro")}
\NormalTok{p\_solo\_bien\_transporte }\OtherTok{\textless{}{-}} \FunctionTok{sample}\NormalTok{(}\DecValTok{4}\SpecialCharTok{:}\DecValTok{10}\NormalTok{, }\DecValTok{1}\NormalTok{)}

\CommentTok{\# Calcular el porcentaje restante para "solo bien\_vivienda1" (equivalente a "solo apartamento")}
\NormalTok{p\_solo\_vivienda1 }\OtherTok{\textless{}{-}}\NormalTok{ total\_puntos\_porcentaje }\SpecialCharTok{{-}}\NormalTok{ (p\_bien\_transporte\_y\_vivienda1 }\SpecialCharTok{+}\NormalTok{ p\_solo\_vivienda2 }\SpecialCharTok{+}
\NormalTok{                                              p\_bien\_transporte\_y\_vivienda2 }\SpecialCharTok{+}\NormalTok{ p\_solo\_bien\_transporte)}

\CommentTok{\# Verificar que el porcentaje restante sea positivo y razonable}
\ControlFlowTok{while}\NormalTok{ (p\_solo\_vivienda1 }\SpecialCharTok{\textless{}} \DecValTok{15}\NormalTok{) \{}
  \CommentTok{\# Reajustar los porcentajes si es necesario}
\NormalTok{  p\_bien\_transporte\_y\_vivienda1 }\OtherTok{\textless{}{-}} \FunctionTok{max}\NormalTok{(p\_bien\_transporte\_y\_vivienda1 }\SpecialCharTok{{-}} \DecValTok{1}\NormalTok{, }\DecValTok{18}\NormalTok{)}
\NormalTok{  p\_solo\_vivienda2 }\OtherTok{\textless{}{-}} \FunctionTok{max}\NormalTok{(p\_solo\_vivienda2 }\SpecialCharTok{{-}} \DecValTok{1}\NormalTok{, }\DecValTok{25}\NormalTok{)}
\NormalTok{  p\_bien\_transporte\_y\_vivienda2 }\OtherTok{\textless{}{-}} \FunctionTok{max}\NormalTok{(p\_bien\_transporte\_y\_vivienda2 }\SpecialCharTok{{-}} \DecValTok{1}\NormalTok{, }\DecValTok{15}\NormalTok{)}

  \CommentTok{\# Recalcular el porcentaje restante}
\NormalTok{  p\_solo\_vivienda1 }\OtherTok{\textless{}{-}}\NormalTok{ total\_puntos\_porcentaje }\SpecialCharTok{{-}}\NormalTok{ (p\_bien\_transporte\_y\_vivienda1 }\SpecialCharTok{+}\NormalTok{ p\_solo\_vivienda2 }\SpecialCharTok{+}
\NormalTok{                                                p\_bien\_transporte\_y\_vivienda2 }\SpecialCharTok{+}\NormalTok{ p\_solo\_bien\_transporte)}
\NormalTok{\}}

\CommentTok{\# Vector con todos los porcentajes}
\NormalTok{porcentajes }\OtherTok{\textless{}{-}} \FunctionTok{c}\NormalTok{(}
\NormalTok{  p\_bien\_transporte\_y\_vivienda1,  }\CommentTok{\# Carro y apartamento}
\NormalTok{  p\_solo\_vivienda1,              }\CommentTok{\# Solo apartamento}
\NormalTok{  p\_solo\_vivienda2,              }\CommentTok{\# Solo casa}
\NormalTok{  p\_solo\_bien\_transporte,        }\CommentTok{\# Solo carro}
\NormalTok{  p\_bien\_transporte\_y\_vivienda2  }\CommentTok{\# Carro y casa}
\NormalTok{)}

\CommentTok{\# Validar que la suma sea exactamente 100}
\ControlFlowTok{if}\NormalTok{ (}\FunctionTok{sum}\NormalTok{(porcentajes) }\SpecialCharTok{!=} \DecValTok{100}\NormalTok{) \{}
  \CommentTok{\# Ajustar el valor más grande para que la suma sea 100}
\NormalTok{  idx\_max }\OtherTok{\textless{}{-}} \FunctionTok{which.max}\NormalTok{(porcentajes)}
\NormalTok{  porcentajes[idx\_max] }\OtherTok{\textless{}{-}}\NormalTok{ porcentajes[idx\_max] }\SpecialCharTok{+}\NormalTok{ (}\DecValTok{100} \SpecialCharTok{{-}} \FunctionTok{sum}\NormalTok{(porcentajes))}
\NormalTok{\}}

\CommentTok{\# Asignar los porcentajes a variables individuales}
\NormalTok{p\_bien\_transporte\_y\_vivienda1 }\OtherTok{\textless{}{-}}\NormalTok{ porcentajes[}\DecValTok{1}\NormalTok{]}
\NormalTok{p\_solo\_vivienda1 }\OtherTok{\textless{}{-}}\NormalTok{ porcentajes[}\DecValTok{2}\NormalTok{]}
\NormalTok{p\_solo\_vivienda2 }\OtherTok{\textless{}{-}}\NormalTok{ porcentajes[}\DecValTok{3}\NormalTok{]}
\NormalTok{p\_solo\_bien\_transporte }\OtherTok{\textless{}{-}}\NormalTok{ porcentajes[}\DecValTok{4}\NormalTok{]}
\NormalTok{p\_bien\_transporte\_y\_vivienda2 }\OtherTok{\textless{}{-}}\NormalTok{ porcentajes[}\DecValTok{5}\NormalTok{]}

\CommentTok{\# Aleatorizar el número total de personas que tienen bien\_transporte y bien\_vivienda1}
\CommentTok{\# (equivalente a "carro y apartamento")}
\NormalTok{total\_bien\_transporte\_y\_vivienda1 }\OtherTok{\textless{}{-}} \FunctionTok{sample}\NormalTok{(}\FunctionTok{c}\NormalTok{(}\DecValTok{48}\NormalTok{, }\DecValTok{60}\NormalTok{, }\DecValTok{72}\NormalTok{, }\DecValTok{100}\NormalTok{, }\DecValTok{120}\NormalTok{, }\DecValTok{144}\NormalTok{, }\DecValTok{150}\NormalTok{, }\DecValTok{180}\NormalTok{, }\DecValTok{200}\NormalTok{, }\DecValTok{250}\NormalTok{), }\DecValTok{1}\NormalTok{)}

\CommentTok{\# Calcular el total de personas en la empresa basado en el porcentaje}
\NormalTok{total\_personas }\OtherTok{\textless{}{-}} \FunctionTok{round}\NormalTok{(total\_bien\_transporte\_y\_vivienda1 }\SpecialCharTok{*} \DecValTok{100} \SpecialCharTok{/}\NormalTok{ p\_bien\_transporte\_y\_vivienda1)}

\CommentTok{\# Calcular el número de personas para cada categoría}
\NormalTok{total\_solo\_vivienda1 }\OtherTok{\textless{}{-}} \FunctionTok{round}\NormalTok{(total\_personas }\SpecialCharTok{*}\NormalTok{ p\_solo\_vivienda1 }\SpecialCharTok{/} \DecValTok{100}\NormalTok{)}
\NormalTok{total\_solo\_vivienda2 }\OtherTok{\textless{}{-}} \FunctionTok{round}\NormalTok{(total\_personas }\SpecialCharTok{*}\NormalTok{ p\_solo\_vivienda2 }\SpecialCharTok{/} \DecValTok{100}\NormalTok{)}
\NormalTok{total\_solo\_bien\_transporte }\OtherTok{\textless{}{-}} \FunctionTok{round}\NormalTok{(total\_personas }\SpecialCharTok{*}\NormalTok{ p\_solo\_bien\_transporte }\SpecialCharTok{/} \DecValTok{100}\NormalTok{)}
\NormalTok{total\_bien\_transporte\_y\_vivienda2 }\OtherTok{\textless{}{-}} \FunctionTok{round}\NormalTok{(total\_personas }\SpecialCharTok{*}\NormalTok{ p\_bien\_transporte\_y\_vivienda2 }\SpecialCharTok{/} \DecValTok{100}\NormalTok{)}

\CommentTok{\# Ajustar el total para que coincida con la suma de todas las categorías}
\NormalTok{total\_calculado }\OtherTok{\textless{}{-}}\NormalTok{ total\_bien\_transporte\_y\_vivienda1 }\SpecialCharTok{+}\NormalTok{ total\_solo\_vivienda1 }\SpecialCharTok{+}\NormalTok{ total\_solo\_vivienda2 }\SpecialCharTok{+}
\NormalTok{                  total\_solo\_bien\_transporte }\SpecialCharTok{+}\NormalTok{ total\_bien\_transporte\_y\_vivienda2}
\ControlFlowTok{if}\NormalTok{ (total\_calculado }\SpecialCharTok{!=}\NormalTok{ total\_personas) \{}
  \CommentTok{\# Ajustar la categoría más grande para mantener el total correcto}
\NormalTok{  diferencia }\OtherTok{\textless{}{-}}\NormalTok{ total\_personas }\SpecialCharTok{{-}}\NormalTok{ total\_calculado}
\NormalTok{  idx\_max }\OtherTok{\textless{}{-}} \FunctionTok{which.max}\NormalTok{(}\FunctionTok{c}\NormalTok{(total\_bien\_transporte\_y\_vivienda1, total\_solo\_vivienda1, total\_solo\_vivienda2,}
\NormalTok{                        total\_solo\_bien\_transporte, total\_bien\_transporte\_y\_vivienda2))}

  \ControlFlowTok{if}\NormalTok{ (idx\_max }\SpecialCharTok{==} \DecValTok{1}\NormalTok{) total\_bien\_transporte\_y\_vivienda1 }\OtherTok{\textless{}{-}}\NormalTok{ total\_bien\_transporte\_y\_vivienda1 }\SpecialCharTok{+}\NormalTok{ diferencia}
  \ControlFlowTok{else} \ControlFlowTok{if}\NormalTok{ (idx\_max }\SpecialCharTok{==} \DecValTok{2}\NormalTok{) total\_solo\_vivienda1 }\OtherTok{\textless{}{-}}\NormalTok{ total\_solo\_vivienda1 }\SpecialCharTok{+}\NormalTok{ diferencia}
  \ControlFlowTok{else} \ControlFlowTok{if}\NormalTok{ (idx\_max }\SpecialCharTok{==} \DecValTok{3}\NormalTok{) total\_solo\_vivienda2 }\OtherTok{\textless{}{-}}\NormalTok{ total\_solo\_vivienda2 }\SpecialCharTok{+}\NormalTok{ diferencia}
  \ControlFlowTok{else} \ControlFlowTok{if}\NormalTok{ (idx\_max }\SpecialCharTok{==} \DecValTok{4}\NormalTok{) total\_solo\_bien\_transporte }\OtherTok{\textless{}{-}}\NormalTok{ total\_solo\_bien\_transporte }\SpecialCharTok{+}\NormalTok{ diferencia}
  \ControlFlowTok{else}\NormalTok{ total\_bien\_transporte\_y\_vivienda2 }\OtherTok{\textless{}{-}}\NormalTok{ total\_bien\_transporte\_y\_vivienda2 }\SpecialCharTok{+}\NormalTok{ diferencia}
\NormalTok{\}}

\CommentTok{\# Verificar nuevamente la consistencia}
\NormalTok{total\_recalculado }\OtherTok{\textless{}{-}}\NormalTok{ total\_bien\_transporte\_y\_vivienda1 }\SpecialCharTok{+}\NormalTok{ total\_solo\_vivienda1 }\SpecialCharTok{+}\NormalTok{ total\_solo\_vivienda2 }\SpecialCharTok{+}
\NormalTok{                    total\_solo\_bien\_transporte }\SpecialCharTok{+}\NormalTok{ total\_bien\_transporte\_y\_vivienda2}
\FunctionTok{stopifnot}\NormalTok{(total\_recalculado }\SpecialCharTok{==}\NormalTok{ total\_personas)}

\CommentTok{\# Generar opciones de respuesta}
\CommentTok{\# La respuesta correcta es el número de personas que poseen solo bien\_vivienda1 (apartamento)}
\NormalTok{respuesta\_correcta }\OtherTok{\textless{}{-}}\NormalTok{ total\_solo\_vivienda1}

\CommentTok{\# Generar distractores plausibles}
\NormalTok{distractor1 }\OtherTok{\textless{}{-}} \FunctionTok{round}\NormalTok{(total\_bien\_transporte\_y\_vivienda1 }\SpecialCharTok{*}\NormalTok{ p\_solo\_vivienda1 }\SpecialCharTok{/}\NormalTok{ p\_bien\_transporte\_y\_vivienda1)}
\NormalTok{distractor2 }\OtherTok{\textless{}{-}} \FunctionTok{round}\NormalTok{(total\_personas }\SpecialCharTok{*}\NormalTok{ p\_solo\_vivienda1 }\SpecialCharTok{/}\NormalTok{ p\_bien\_transporte\_y\_vivienda1)}
\NormalTok{distractor3 }\OtherTok{\textless{}{-}}\NormalTok{ total\_personas }\SpecialCharTok{{-}}\NormalTok{ total\_solo\_vivienda1}

\CommentTok{\# Asegurar que los distractores sean diferentes entre sí y de la respuesta correcta}
\ControlFlowTok{while}\NormalTok{ (}\FunctionTok{length}\NormalTok{(}\FunctionTok{unique}\NormalTok{(}\FunctionTok{c}\NormalTok{(respuesta\_correcta, distractor1, distractor2, distractor3))) }\SpecialCharTok{\textless{}} \DecValTok{4}\NormalTok{) \{}
\NormalTok{  distractor1 }\OtherTok{\textless{}{-}}\NormalTok{ distractor1 }\SpecialCharTok{+} \FunctionTok{sample}\NormalTok{(}\SpecialCharTok{{-}}\DecValTok{5}\SpecialCharTok{:}\DecValTok{5}\NormalTok{, }\DecValTok{1}\NormalTok{)}
\NormalTok{  distractor2 }\OtherTok{\textless{}{-}}\NormalTok{ distractor2 }\SpecialCharTok{+} \FunctionTok{sample}\NormalTok{(}\SpecialCharTok{{-}}\DecValTok{5}\SpecialCharTok{:}\DecValTok{5}\NormalTok{, }\DecValTok{1}\NormalTok{)}
\NormalTok{  distractor3 }\OtherTok{\textless{}{-}}\NormalTok{ distractor3 }\SpecialCharTok{+} \FunctionTok{sample}\NormalTok{(}\SpecialCharTok{{-}}\DecValTok{5}\SpecialCharTok{:}\DecValTok{5}\NormalTok{, }\DecValTok{1}\NormalTok{)}
\NormalTok{\}}

\CommentTok{\# Crear el vector de opciones y mezclarlas}
\NormalTok{opciones }\OtherTok{\textless{}{-}} \FunctionTok{c}\NormalTok{(respuesta\_correcta, distractor1, distractor2, distractor3)}
\NormalTok{indice\_correcto }\OtherTok{\textless{}{-}} \DecValTok{1}  \CommentTok{\# Posición actual de la respuesta correcta}
\NormalTok{opciones\_mezcladas }\OtherTok{\textless{}{-}} \FunctionTok{sample}\NormalTok{(opciones)}
\NormalTok{indice\_correcto }\OtherTok{\textless{}{-}} \FunctionTok{which}\NormalTok{(opciones\_mezcladas }\SpecialCharTok{==}\NormalTok{ respuesta\_correcta)}

\CommentTok{\# Vector de solución para r{-}exams (0 = falso, 1 = verdadero)}
\NormalTok{solucion }\OtherTok{\textless{}{-}} \FunctionTok{rep}\NormalTok{(}\DecValTok{0}\NormalTok{, }\DecValTok{4}\NormalTok{)}
\NormalTok{solucion[indice\_correcto] }\OtherTok{\textless{}{-}} \DecValTok{1}

\CommentTok{\# Aleatorizar colores para el gráfico circular usando colores compatibles con matplotlib}
\CommentTok{\# Definir paletas de colores para matplotlib (nombres de colores estándar y códigos hexadecimales)}
\NormalTok{paleta\_colores1 }\OtherTok{\textless{}{-}} \FunctionTok{c}\NormalTok{(}\StringTok{"\#4C72B0"}\NormalTok{, }\StringTok{"\#55A868"}\NormalTok{, }\StringTok{"\#C44E52"}\NormalTok{, }\StringTok{"\#8172B3"}\NormalTok{, }\StringTok{"\#CCB974"}\NormalTok{, }\StringTok{"\#64B5CD"}\NormalTok{)}
\NormalTok{paleta\_colores2 }\OtherTok{\textless{}{-}} \FunctionTok{c}\NormalTok{(}\StringTok{"\#1F77B4"}\NormalTok{, }\StringTok{"\#FF7F0E"}\NormalTok{, }\StringTok{"\#2CA02C"}\NormalTok{, }\StringTok{"\#D62728"}\NormalTok{, }\StringTok{"\#9467BD"}\NormalTok{, }\StringTok{"\#8C564B"}\NormalTok{)}
\NormalTok{paleta\_colores3 }\OtherTok{\textless{}{-}} \FunctionTok{c}\NormalTok{(}\StringTok{"\#E41A1C"}\NormalTok{, }\StringTok{"\#377EB8"}\NormalTok{, }\StringTok{"\#4DAF4A"}\NormalTok{, }\StringTok{"\#984EA3"}\NormalTok{, }\StringTok{"\#FF7F00"}\NormalTok{, }\StringTok{"\#FFFF33"}\NormalTok{)}

\CommentTok{\# Seleccionar una paleta aleatoria}
\NormalTok{paleta\_seleccionada }\OtherTok{\textless{}{-}} \FunctionTok{sample}\NormalTok{(}\FunctionTok{list}\NormalTok{(paleta\_colores1, paleta\_colores2, paleta\_colores3), }\DecValTok{1}\NormalTok{)[[}\DecValTok{1}\NormalTok{]]}

\CommentTok{\# Mezclar la paleta para obtener combinaciones aleatorias}
\NormalTok{paleta\_mezclada }\OtherTok{\textless{}{-}} \FunctionTok{sample}\NormalTok{(paleta\_seleccionada)}

\CommentTok{\# Asignar colores a cada sección del gráfico}
\NormalTok{color\_bien\_transporte\_y\_vivienda1 }\OtherTok{\textless{}{-}}\NormalTok{ paleta\_mezclada[}\DecValTok{1}\NormalTok{]}
\NormalTok{color\_solo\_vivienda1 }\OtherTok{\textless{}{-}}\NormalTok{ paleta\_mezclada[}\DecValTok{2}\NormalTok{]}
\NormalTok{color\_solo\_vivienda2 }\OtherTok{\textless{}{-}}\NormalTok{ paleta\_mezclada[}\DecValTok{3}\NormalTok{]}
\NormalTok{color\_solo\_bien\_transporte }\OtherTok{\textless{}{-}}\NormalTok{ paleta\_mezclada[}\DecValTok{4}\NormalTok{]}
\NormalTok{color\_bien\_transporte\_y\_vivienda2 }\OtherTok{\textless{}{-}}\NormalTok{ paleta\_mezclada[}\DecValTok{5}\NormalTok{]}
\end{Highlighting}
\end{Shaded}

\begin{Shaded}
\begin{Highlighting}[]
\FunctionTok{options}\NormalTok{(}\AttributeTok{OutDec =} \StringTok{"."}\NormalTok{)  }\CommentTok{\# Asegurar punto decimal en este chunk}

\CommentTok{\# Código Python para generar el gráfico circular}
\NormalTok{codigo\_python\_grafico }\OtherTok{\textless{}{-}} \FunctionTok{sprintf}\NormalTok{(}
\StringTok{\textquotesingle{}}
\StringTok{import matplotlib.pyplot as plt}
\StringTok{import numpy as np}

\StringTok{\# Datos para el gráfico circular}
\StringTok{categorias = ["\%s y \%s", "Solo \%s", "Solo \%s", "Solo \%s", "\%s y \%s"]}
\StringTok{porcentajes = [\%d, \%d, \%d, \%d, \%d]}

\StringTok{\# Crear los valores exactos como etiquetas}
\StringTok{etiquetas = [f"\{categorias[0]\}}\SpecialCharTok{\textbackslash{}\textbackslash{}}\StringTok{n\{porcentajes[0]\}\%\%",}
\StringTok{            f"\{categorias[1]\}}\SpecialCharTok{\textbackslash{}\textbackslash{}}\StringTok{n\{porcentajes[1]\}\%\%",}
\StringTok{            f"\{categorias[2]\}}\SpecialCharTok{\textbackslash{}\textbackslash{}}\StringTok{n\{porcentajes[2]\}\%\%",}
\StringTok{            f"\{categorias[3]\}}\SpecialCharTok{\textbackslash{}\textbackslash{}}\StringTok{n\{porcentajes[3]\}\%\%",}
\StringTok{            f"\{categorias[4]\}}\SpecialCharTok{\textbackslash{}\textbackslash{}}\StringTok{n\{porcentajes[4]\}\%\%"]}

\StringTok{\# Aleatorizar posición de inicio (ángulo)}
\StringTok{angulo\_inicio = np.random.randint(0, 360)}
\StringTok{angulo\_separacion = 0.05  \# Separación entre segmentos}

\StringTok{\# Configuración de colores}
\StringTok{colores = ["\%s", "\%s", "\%s", "\%s", "\%s"]}

\StringTok{\# Crear el gráfico circular}
\StringTok{plt.figure(figsize=(6, 5))}
\StringTok{cunas, textos = plt.pie(}
\StringTok{    porcentajes,}
\StringTok{    labels=None,}
\StringTok{    colors=colores,}
\StringTok{    startangle=angulo\_inicio,}
\StringTok{    explode=[0.05, 0.05, 0.05, 0.05, 0.05],  \# Separación de todos los segmentos}
\StringTok{    wedgeprops=\{"edgecolor": "white", "linewidth": 1\}}
\StringTok{)}

\StringTok{\# Añadir etiquetas con porcentajes fuera del gráfico}
\StringTok{plt.legend(cunas, etiquetas, loc="center left", bbox\_to\_anchor=(1, 0.5))}

\StringTok{\# Ajustar diseño}
\StringTok{plt.axis("equal")}
\StringTok{plt.tight\_layout()}

\StringTok{\# Guardar el gráfico}
\StringTok{plt.savefig("grafico\_circular.png", dpi=150, bbox\_inches="tight")}
\StringTok{plt.close()}
\StringTok{\textquotesingle{}}\NormalTok{,}
\NormalTok{bien\_transporte, bien\_vivienda1, bien\_vivienda1, bien\_vivienda2, bien\_transporte, bien\_transporte, bien\_vivienda2,}
\NormalTok{p\_bien\_transporte\_y\_vivienda1, p\_solo\_vivienda1, p\_solo\_vivienda2, p\_solo\_bien\_transporte, p\_bien\_transporte\_y\_vivienda2,}
\NormalTok{color\_bien\_transporte\_y\_vivienda1, color\_solo\_vivienda1, color\_solo\_vivienda2, color\_solo\_bien\_transporte, color\_bien\_transporte\_y\_vivienda2}
\NormalTok{)}

\CommentTok{\# Ejecutar el código Python}
\FunctionTok{py\_run\_string}\NormalTok{(codigo\_python\_grafico)}
\end{Highlighting}
\end{Shaded}

\section{Question}\label{question}

Se realizó un(a) censo a un grupo de personas de una organización de
servicios sobre el tipo de bienes que poseen. Los resultados se
presentan en la gráfica.

\begin{figure}

{\centering \includegraphics[width=.75\linewidth]{grafico_circular} 

}

\end{figure}

Si 100 personas de el(la) organización poseen automóvil y(e) piso,
¿cuántas personas poseen solo piso?

\subsection{Answerlist}\label{answerlist}

\begin{itemize}
\tightlist
\item
  70
\item
  261
\item
  66
\item
  305
\end{itemize}

\section{Solution}\label{solution}

Para resolver este problema, debemos usar las relaciones proporcionadas
en el gráfico circular y aplicar proporciones.

Primero, identificamos los datos conocidos:

\begin{itemize}
\tightlist
\item
  27\% de las personas tienen automóvil y piso
\item
  19\% de las personas tienen solo piso (este es el dato que buscamos)
\item
  Sabemos que 100 personas tienen automóvil y piso
\end{itemize}

Paso 1: Calcular el número total de personas en la organización. Si 100
personas representan el 27\% del total, entonces:

\[\text{Total de personas} = \frac{\text{Personas con automóvil y piso} \times 100\%}{\text{Porcentaje de personas con automóvil y piso}}\]

\[\text{Total de personas} = \frac{100 \times 100\%}{27\%} = 370\]

Paso 2: Calcular cuántas personas tienen solo piso. El porcentaje de
personas que tienen solo piso es 19\% del total, por lo tanto:

\[\text{Personas con solo piso} = \text{Total de personas} \times \frac{\text{Porcentaje de personas con solo piso}}{100\%}\]

\[\text{Personas con solo piso} = 370 \times \frac{19\%}{100\%} = 70\]

Por lo tanto, 70 personas poseen solo piso.

\subsection{Answerlist}\label{answerlist-1}

\begin{itemize}
\tightlist
\item
  Verdadero
\item
  Falso
\item
  Falso
\item
  Falso
\end{itemize}

\section{Meta-information}\label{meta-information}

exname: interpretacion\_grafico\_circular\_proporciones extype: schoice
exsolution: 1000 exshuffle: TRUE exsection: Interpretación de
gráficos\textbar Proporciones\textbar Porcentajes extol: 0
exextra{[}tipoicfes{]}: SAI exextra{[}nivel{]}: Básico
exextra{[}componente{]}: Aleatorio exextra{[}competencia{]}:
Interpretación y representación exextra{[}tema{]}: Interpretación de
gráficos circulares

\end{document}
